\chapter{Complexity — P Subset P Poly Config Circuit}
%%% generated by the blueprint pipeline from TCSlib.Complexity.CircuitComplexity.PSubsetPPolyConfigCircuit
\section{Overview}

Declarations of the Lean module \texttt{TCSlib.Complexity.CircuitComplexity.PSubsetPPolyConfigCircuit}.

\section{Declarations}

\begin{theorem}[Indexing a fixed-length flat map chunk]
\lean{Complexity.getD_flatMap_chunk}
\label{Complexity.getD_flatMap_chunk}
\leanok
\uses{Complexity.CfgTableau.take_drop_flatMap_append}
\difficulty{3}
Let $l$ be a finite list over a type $\alpha$, let $g$ map each element of $\alpha$ to a
list over $\beta$, and suppose every list $g(a)$ has length $c$. For any $i<|l|$ and
$j<c$, and any default value $d\in\beta$, the element at index $ci+j$ of the
concatenation of the lists $g(a)$ for $a$ in $l$, followed by a list $rest$, with
default $d$, equals the element at index $j$ of $g(l_i)$ with default $d$.
\end{theorem}

\begin{theorem}[Correctness of the input-cell update]
\lean{Complexity.inp_block}
\label{Complexity.inp_block}
\leanok
\uses{BoolCircuit.BitSrc, BoolCircuit.BitSrc.eval, Complexity.CfgTableau.cfgAt, Complexity.CfgTableau.optBits, Complexity.cfgExp, Complexity.cfgExp_inpIdx, Complexity.cfgF, Complexity.cfgF_inp_succ, Complexity.cfgF_inp_zero, Complexity.drop_take_padSrcs, Complexity.eval_block, Complexity.ichSrc, Complexity.ihSrc, Complexity.inpChain, Complexity.inpExp, Complexity.inpIdx, Complexity.inpIdx_lt_succ, Complexity.inpSrcs, Complexity.inputBitAt, Complexity.length_prevSnapSrcs, Complexity.padSrcs, Complexity.prevSnapSrcs, Complexity.prevSnapSrcs_eval, Complexity.snapDecode_snapEncode, Complexity.snapIdx_lt_succ}
\difficulty{7}
Let $M$ be a finite-state multi-tape Turing machine over the binary alphabet, let $\ell$
be a list of circuit bit sources, let $T,t,p\in\mathbb N$, let $x$ be a binary string,
and let $B$ be a list of binary blocks. Suppose every source $s$ in $\ell$ has the same
value on $(x,B)$ as on $(x,[])$. Suppose also that every block of index
$i'<t\bigl(k(2T+1)+|\ell|+3\bigr)+k(2T+1)+p$ equals the configuration encoding at index
$i'$ generated from the list of source values on $(x,[])$, where $k$ is the number of
work tapes of $M$. If $p<|\ell|+2$, then applying the input-cell update to the values of
the padded input-cell source list for time $t$ and position $p$ gives exactly the
input-cell encoding at that time and position for the binary input whose bits are the
values of the sources in $\ell$ on $(x,[])$.
\end{theorem}

\begin{definition}[Snapshot work sources]
\lean{Complexity.snapWorkSrcs}
\label{Complexity.snapWorkSrcs}
\leanok
\uses{BoolCircuit.BitSrc, Complexity.cellIdx}
Let $M$ be a finite Turing machine with Boolean symbols, let $\ell$ be a finite list of
bit sources, and let $T,t\in\mathbb{N}$. For each machine state $\tau$, in order, take
bits $3$ and $4$ of the cell block at time $t$, state $\tau$, and position $2T$;
concatenate these pairs to form the snapshot work-source list.
\end{definition}

\begin{theorem}[Length of work-tape snapshot sources]
\lean{Complexity.length_snapWorkSrcs}
\label{Complexity.length_snapWorkSrcs}
\leanok
\uses{Complexity.CfgTableau.length_flatMap_const, Complexity.snapWorkSrcs}
\difficulty{3}
Let $M$ be a finite-state Turing machine over the binary alphabet, let $\ell$ be a list
of bit sources, and let $T,t$ be natural numbers. The list of bit sources encoding the
work-tape contents in the snapshot at time $t$ has length $2k$, where $k$ is the number
of work tapes of $M$.
\end{theorem}

\begin{theorem}[Correctness of a snapshot block]
\lean{Complexity.snap_block}
\label{Complexity.snap_block}
\leanok
\uses{BoolCircuit.BitSrc, BoolCircuit.BitSrc.eval, Complexity.CfgTableau.cfgAt, Complexity.CfgTableau.optBits, Complexity.cellExp, Complexity.cellIdx, Complexity.cellIdx_lt_inpIdx, Complexity.cfgExp, Complexity.cfgExp_cellIdx, Complexity.cfgExp_inpIdx, Complexity.cfgExp_snapIdx, Complexity.cfgF, Complexity.cfgF_snap, Complexity.drop_take_padSrcs, Complexity.eval_block, Complexity.getD_flatMap_chunk, Complexity.inpChain_right, Complexity.inpExp, Complexity.inpIdx, Complexity.inpIdx_lt_snapIdx, Complexity.length_prevSnapSrcs, Complexity.length_snapWorkSrcs, Complexity.padSrcs, Complexity.prevSnapSrcs, Complexity.prevSnapSrcs_eval, Complexity.snapDecode_snapEncode, Complexity.snapExp, Complexity.snapIdx, Complexity.snapIdx_lt_succ, Complexity.snapSrcs, Complexity.snapWidth, Complexity.snapWorkSrcs, Complexity.workChain_right, Turing.Cfg.inputSymbol, Turing.Cfg.workTapeSymbols}
\difficulty{6}
Let $M$ be a finite binary Turing machine with $k$ work tapes, let $\ell$ be a finite
list of bit sources, and let $T,t\in\mathbb{N}$ with $t\le T$. For every input word $x$
and list of Boolean blocks such that, at each index $i$ below the snapshot-block index
for time $t$, block $i$ equals the configuration expansion specified by $\ell$ on the
source values from $x$, evaluating the padded snapshot sources for time $t$ on $x$ and
those blocks and applying the configuration update for a snapshot gives exactly the
snapshot expansion of $M$ on the source values from $x$ at time $t$.
\end{theorem}

\begin{theorem}[Strict Bound from a Common Lower Bound]
\lean{Complexity.prev_lt}
\label{Complexity.prev_lt}
\leanok
\uses{Complexity.cfgStride}
\difficulty{1}
Let $M$ be a finite-state Turing machine with Boolean tape alphabet, let $\ell$ be a
list of bit sources, and let $T,t,i,a$ be natural numbers. Define
$S=M.k(2T+1)+|\ell|+3$. If $a<tS$ and $tS\le i$, then $a<i$.
\end{theorem}

\begin{theorem}[Previous Configuration Index Bound]
\lean{Complexity.prevStep_lt}
\label{Complexity.prevStep_lt}
\leanok
\uses{Complexity.cfgStride}
\difficulty{3}
Let $M$ be a finite-state Turing machine, let $\ell$ be a finite list of bit sources,
and let $T,t,a,i\in\mathbb N$. If $t\ne0$, $a<t\,s$, and $t\,s\leq i$, where
$s=k(2T+1)+|\ell|+3$ and $k$ is the number of work tapes of $M$, then $a<i$.
\end{theorem}

\begin{theorem}[Validity of previous-snapshot sources]
\lean{Complexity.prevSnapSrcs_valid}
\label{Complexity.prevSnapSrcs_valid}
\leanok
\uses{BoolCircuit.BitSrc.Valid, Complexity.cfgStride, Complexity.cfgWidth, Complexity.prevSnapSrcs, Complexity.prevStep_lt, Complexity.snapIdx_lt_succ}
\difficulty{4}
Let $M$ be a finite Turing machine with $k$ work tapes, let $\ell$ be a finite list of
bit sources, and let $T,n,t,i$ be natural numbers. If
$t\,\mathrm{cfgStride}(M,|\ell|,T)\le i$, then every source in the list encoding the
previous snapshot at time $t$ is valid for input width $n$, configuration width
$\mathrm{cfgWidth}(M)$, and instruction index $i$.
\end{theorem}

\begin{theorem}[Validity of cell sources]
\lean{Complexity.cellSrcs_valid}
\label{Complexity.cellSrcs_valid}
\leanok
\uses{BoolCircuit.BitSrc.Valid, Complexity.cellIdx, Complexity.cellIdx_lt_succ, Complexity.cellSrcs, Complexity.cfgWidth, Complexity.chSrc, Complexity.le_cellIdx, Complexity.nbSrc, Complexity.prevSnapSrcs_valid, Complexity.prevStep_lt}
\difficulty{5}
Let $M$ be a finite Turing machine over the binary alphabet, let $\ell$ be a finite list
of bit sources, and let $T,n,t\in\mathbb{N}$, $\tau<k$, and $r<2T+1$, where $k$ is the
number of work tapes of $M$. Every source in the list constructed for cell $(t,\tau,r)$
is valid with input length $n$, block width equal to the snapshot width of $M$ plus $5$,
and current block index equal to the cell index of $(t,\tau,r)$; that is, each input
source refers to an index below $n$, and each block source refers to an earlier block
and a bit index below that width.
\end{theorem}

\begin{theorem}[Validity of input-cell sources]
\lean{Complexity.inpSrcs_valid}
\label{Complexity.inpSrcs_valid}
\leanok
\uses{BoolCircuit.BitSrc.Valid, Complexity.cfgStride, Complexity.cfgWidth, Complexity.ichSrc, Complexity.ihSrc, Complexity.inpIdx, Complexity.inpIdx_lt_succ, Complexity.inpSrcs, Complexity.prevSnapSrcs_valid, Complexity.prevStep_lt}
\difficulty{5}
Let $M$ be a finite-state Turing machine with $k$ work tapes, and let
$n,T\in\mathbb{N}$. Let $\ell$ be a list of bit sources, each valid for input length
$n$, width $W=\operatorname{snapWidth}(M)+5$, and block index $0$. For natural numbers
$t,p$ with $p<|\ell|+2$, every source in the input-cell source list at time $t$ and
position $p$ is valid for input length $n$, width $W$, and block index
$t\bigl(k(2T+1)+|\ell|+3\bigr)+k(2T+1)+|\ell|+2+p$.
\end{theorem}

\begin{theorem}[Validity of snapshot sources]
\lean{Complexity.snapSrcs_valid}
\label{Complexity.snapSrcs_valid}
\leanok
\uses{BoolCircuit.BitSrc.Valid, Complexity.cellIdx_lt_inpIdx, Complexity.cfgStride, Complexity.cfgWidth, Complexity.inpIdx_lt_snapIdx, Complexity.prevSnapSrcs_valid, Complexity.prevStep_lt, Complexity.snapIdx, Complexity.snapIdx_lt_succ, Complexity.snapSrcs}
\difficulty{5}
Let $M$ be a finite Turing machine over the Boolean alphabet, let $\ell$ be a finite
list of bit sources, and let $T,t,n$ be natural numbers. Every source in the
snapshot-source list for $M$, $\ell$, and time $t$ is valid for input width $n$,
configuration width equal to the number of possible snapshots of $M$ plus five, and
block index corresponding to time $t$.
\end{theorem}

\begin{theorem}[Length of padded bit sources]
\lean{Complexity.length_padSrcs}
\label{Complexity.length_padSrcs}
\leanok
\uses{BoolCircuit.BitSrc, Complexity.cfgArity, Complexity.padSrcs}
\difficulty{2}
Let $M$ be a finite-state Turing machine over the binary alphabet with $k$ work tapes,
and let $l$ be a finite list of bit sources. If $|l|$ is at most the configuration arity
of $M$ (the number of possible snapshots of $M$, plus $2k+13$), then padding $l$ with
false constants to that arity produces a list of length exactly the configuration arity.
\end{theorem}

\begin{theorem}[Configuration program well-formedness]
\lean{Complexity.cfgProg_wf}
\label{Complexity.cfgProg_wf}
\leanok
\uses{BoolCircuit.BitSrc, BoolCircuit.BitSrc.Valid, BoolCircuit.ProgWF, Complexity.cellSrcs, Complexity.cellSrcs_valid, Complexity.cfgArity, Complexity.cfgIdx_cases, Complexity.cfgInstr, Complexity.cfgPos_cellIdx, Complexity.cfgPos_inpIdx, Complexity.cfgPos_snapIdx, Complexity.cfgProg, Complexity.cfgProg_getElem, Complexity.cfgWidth, Complexity.inpSrcs, Complexity.inpSrcs_valid, Complexity.length_cfgProg, Complexity.length_padSrcs, Complexity.padSrcs, Complexity.prevSnapSrcs, Complexity.snapSrcs, Complexity.snapSrcs_valid, Complexity.snapWorkSrcs}
\difficulty{6}
Let $M$ be a finite-state Turing machine with $k$ work tapes, and let $n,T$ be natural
numbers. Suppose every source in the input list $\ell$ is either an input bit with index
less than $n$, a Boolean constant, or a bit of a block with block index less than $0$
and bit index less than the configuration width. Then every source used by an
instruction at index $i$ in the configuration program is an input bit of index less than
$n$, a constant, or a bit of a block with index less than $i$ and bit index less than
the configuration width; moreover, every instruction in the program has source-list
length equal to the configuration arity.
\end{theorem}

\begin{theorem}[Evaluation of a valid source at layout zero]
\lean{Complexity.eval_layout}
\label{Complexity.eval_layout}
\leanok
\uses{BoolCircuit.BitSrc, BoolCircuit.BitSrc.Valid, BoolCircuit.BitSrc.eval, Complexity.cfgWidth}
\difficulty{3}
Let $M$ be a finite Turing machine over the Boolean alphabet, let $n$ be a natural
number, and let $s$ be a bit source valid for input width $n$, block width equal to the
number of possible snapshots of $M$ plus $5$, and layout index $0$. For any list $x$ of
input bits and any list $bs$ of blocks of bits, evaluating $s$ on $x$ and $bs$ gives the
same Boolean value as evaluating $s$ on $x$ and an empty list of blocks.
\end{theorem}

\begin{theorem}[Configuration program blocks]
\lean{Complexity.cfgProg_blocks}
\label{Complexity.cfgProg_blocks}
\leanok
\uses{BoolCircuit.BitSrc.Valid, BoolCircuit.BitSrc.eval, BoolCircuit.progBlocks, BoolCircuit.progBlocks_getD, Complexity.CfgTableau.cell_block, Complexity.cfgExp, Complexity.cfgExp_cellIdx, Complexity.cfgExp_inpIdx, Complexity.cfgExp_snapIdx, Complexity.cfgF, Complexity.cfgIdx_cases, Complexity.cfgInstr, Complexity.cfgPos_cellIdx, Complexity.cfgPos_inpIdx, Complexity.cfgPos_snapIdx, Complexity.cfgProg, Complexity.cfgProg_getElem, Complexity.cfgProg_wf, Complexity.cfgWidth, Complexity.eval_layout, Complexity.inp_block, Complexity.length_cfgProg, Complexity.snap_block}
\difficulty{5}
Let $M$ be a finite-state Turing machine with $k$ work tapes, and let $n,T$ be natural
numbers. Let $\ell$ be a list of bit sources, each of which is either an input bit with
index less than $n$, a Boolean constant, or a bit of a block with index less than $0$
and bit index less than the configuration width; let $x$ be any Boolean list. For every
index $i$ within the configuration program, its computed block is the configuration
expansion at index $i$, evaluated on the input obtained by evaluating each source in
$\ell$ against $x$ and the empty list of prior blocks.
\end{theorem}

\begin{definition}[Configuration circuit]
\lean{Complexity.cfgCircuit}
\label{Complexity.cfgCircuit}
\leanok
\uses{BoolCircuit.BitSrc, BoolCircuit.DAGCircuit, BoolCircuit.progCircuit, Complexity.cfgArity, Complexity.cfgF, Complexity.cfgProg, Complexity.cfgWidth, Complexity.snapIdx, Complexity.snapWidth}
Let $M$ be a finite Turing machine with Boolean symbols, let $n,T\in\mathbb{N}$, and let
$\ell$ be a finite list of bit sources. The configuration circuit on $n$ input bits is
the directed acyclic circuit obtained from the configuration program using the machine's
configuration arity, width, and Boolean operation, with output vertex the snapshot block
at time $T$ and output width the snapshot width.
\end{definition}

\begin{definition}[Configuration constant]
\lean{Complexity.cfgConst}
\label{Complexity.cfgConst}
\leanok
\uses{BoolCircuit.kindWidth, Complexity.cfgArity, Complexity.cfgF, Complexity.cfgWidth}
For a finite binary Turing machine $M$, define its configuration constant to be the
supremum, over all configuration kinds, of the width bound evaluated at $M$’s
configuration arity, configuration width, and local update function.
\end{definition}

\begin{theorem}[Evaluation of the configuration circuit]
\lean{Complexity.cfgCircuit_eval}
\label{Complexity.cfgCircuit_eval}
\leanok
\uses{BoolCircuit.BitSrc.Valid, BoolCircuit.BitSrc.eval, BoolCircuit.DAGCircuit.eval, BoolCircuit.progCircuit_eval, Complexity.cfgCircuit, Complexity.cfgExp_snapIdx, Complexity.cfgProg, Complexity.cfgProg_blocks, Complexity.cfgProg_wf, Complexity.cfgWidth, Complexity.emitted, Complexity.length_cfgProg, Complexity.snapExp, Complexity.snapIdx, Complexity.snapIdx_lt_succ, Complexity.snapshotAt, Complexity.tableauAccepted}
\difficulty{3}
Let $M$ be a finite-state Turing machine with $k$ work tapes and binary input and
output, let $n,T\in\mathbb{N}$, let $\ell$ be a finite list of bit sources valid for $n$
input bits, and let $x\in\{0,1\}^n$. Form the binary word $y$ by evaluating each source
in $\ell$ on $x$. The configuration circuit built from $M, n,\ell,T$ evaluates to true
exactly when, at some time $s<T$ in the run of $M$ on $y$, the transition from the
machine's snapshot emits the bit $1$.
\end{theorem}

\begin{theorem}[Fan-in bound for the configuration circuit]
\lean{Complexity.cfgCircuit_isFaninTwo}
\label{Complexity.cfgCircuit_isFaninTwo}
\leanok
\uses{BoolCircuit.DAGCircuit.IsFaninTwo, BoolCircuit.progCircuit_isFaninTwo, Complexity.cfgCircuit}
\difficulty{1}
For every finite-state binary Turing machine $M$, natural number $n$, finite list of bit
sources $\ell$, and time bound $T$, the configuration circuit built from $M,n,\ell,T$ is
well formed and every gate reads at most two vertices. Here well formed means that each
gate has no repeated argument and every NOT gate has exactly one argument.
\end{theorem}

\begin{theorem}[Size bound for the configuration circuit]
\lean{Complexity.cfgCircuit_size_le}
\label{Complexity.cfgCircuit_size_le}
\leanok
\uses{BoolCircuit.DAGCircuit.size, BoolCircuit.progCircuit_size_le, Complexity.cfgArity, Complexity.cfgCircuit, Complexity.cfgConst, Complexity.cfgF, Complexity.cfgProg, Complexity.cfgWidth, Complexity.length_cfgProg, Complexity.snapIdx, Complexity.snapWidth}
\difficulty{3}
Let $M$ be a finite-state binary Turing machine with $k$ work tapes, let
$n,T\in\mathbb{N}$, and let $\ell$ be a finite list of bit sources. Let $K$ be the
maximum, over the finitely many instruction kinds of the configuration computation for
$M$, of the total width of the circuits used to compute that instruction's output bits.
Then the size of the configuration circuit is at most
$n+2+(T+1)\bigl(k(2T+1)+|\ell|+3\bigr)K$.
\end{theorem}
